\section{Taksonomi Kesalahan Program dan Alat Debugging}
\label{sec:taksonomi-kesalahan}

Menulis program assembly berarti berhadapan langsung dengan register,
\eng{flag}, dan alamat memori tanpa lapisan pelindung apa pun. Satu
digit salah pada operand, satu instruksi yang tertukar, atau satu
\asm{PUSH} yang lupa dipasangkan dengan \asm{POP} dapat membuat
program berhenti bekerja. Karena itu kemampuan menemukan kesalahan
menjadi keterampilan yang setara pentingnya dengan kemampuan menulis
program itu sendiri.
\begin{definisi}
\term{Debugging} adalah proses sistematis untuk menemukan penyebab
suatu kesalahan (\eng{bug}) pada program, membuktikan penyebab itu
melalui bukti pengamatan, lalu memperbaiki dan memverifikasi hasil
perbaikan. \eng{Debugging} bukan kegiatan menebak-nebak, melainkan
kegiatan mengumpulkan bukti pada level instruksi.
\end{definisi}

Pada bab-bab sebelumnya telah dibahas bagaimana register bekerja
(Bab~\ref{chap:register}), bagaimana keputusan diambil berdasarkan
\eng{flag} (Bab~\ref{chap:kontrol}), dan bagaimana prosedur
menggunakan \eng{stack} (Bab~\ref{chap:prosedur}). Ketiga pengetahuan
itu menjadi bekal utama pada bab ini.

\subsection{Lima Jenis Kesalahan Program}

Tidak semua kesalahan sama. Mengenali jenis kesalahan lebih dahulu
mempercepat pemilihan alat dan strategi. Secara umum terdapat lima
jenis kesalahan yang perlu dibedakan.

\begin{description}
    \item[Kesalahan sintaks] Kesalahan penulisan yang membuat
          assembler menolak berkas sumber sehingga program
          \emph{tidak} dihasilkan sama sekali. Contohnya mnemonic yang
          salah tulis (\code{MOVV} alih-alih \asm{MOV}), operand yang
          ukurannya tidak cocok, atau label yang tidak pernah
          didefinisikan. Jenis ini paling mudah diperbaiki karena
          assembler menunjukkan nomor barisnya.

    \item[Kesalahan saat berjalan] Kesalahan yang baru muncul ketika
          program dijalankan, misalnya pembagian dengan nol
          (\asm{DIV} dengan pembagi nol) atau pembacaan berkas yang
          tidak ada. Program berjalan sampai titik tertentu lalu
          berhenti dengan pesan dari sistem operasi.

    \item[Kesalahan logika] Kesalahan yang paling sulit: program
          \emph{berjalan sampai selesai} tetapi hasilnya salah.
          Assembler diam, sistem operasi pun diam. Contohnya
          penghitung perulangan yang salah nilai awal, operasi
          \asm{ADD} yang seharusnya \asm{SUB}, atau akumulasi yang
          memakai register yang keliru.

    \item[Kesalahan segmentasi] Kesalahan karena program mengakses
          alamat di luar segmen yang seharusnya, misalnya karena
          register segmen belum diinisialisasi, indeks larik melewati
          batas, atau \reg{DS} tidak pernah diarahkan ke segmen data.

    \item[\eng{Overflow}] Kesalahan karena hasil perhitungan atau
          pemakaian \eng{stack} melebihi daya tampungnya.
          \eng{Overflow} aritmatika terjadi ketika hasil melewati
          batas register 16-bit, misalnya lebih dari \hex{FFFF}.
          \eng{Overflow} \eng{stack} terjadi ketika \asm{PUSH} atau
          \asm{CALL} melebihi ruang \dir{.STACK} yang disediakan,
          atau ketika jumlah \asm{PUSH} dan \asm{POP} tidak seimbang
          sehingga \asm{RET} mengambil alamat yang salah.
\end{description}

\begin{table}[H]
    \centering
    \caption{Taksonomi kesalahan program beserta gejala dan cara
    penemuannya.}
    \label{tab:taksonomi-kesalahan}
    \small
    \begin{tabular}{@{}p{2.1cm}p{3.6cm}p{3.3cm}p{4.0cm}@{}}
        \toprule
        \textbf{Jenis} & \textbf{Gejala} & \textbf{Kapan
        terdeteksi} & \textbf{Cara menemukan} \\
        \midrule
        Sintaks & Assembler menolak berkas; tidak ada
        \file{.obj} & Sebelum program dijalankan &
        Baca pesan assembler beserta nomor barisnya \\
        \addlinespace
        Saat berjalan & Program berhenti di tengah jalan; pesan
        sistem operasi & Ketika dijalankan &
        Jalankan bertahap sampai titik berhenti \\
        \addlinespace
        Logika & Program selesai normal, hasil salah & Tidak
        terdeteksi otomatis & Tabel \eng{trace} register dan
        \eng{flag} per langkah \\
        \addlinespace
        Segmentasi & Karakter sampah, nilai berubah sendiri,
        program macet & Ketika dijalankan &
        Periksa \reg{DS}, \reg{ES}, \reg{SS}, dan batas indeks
        melalui \code{D} \\
        \addlinespace
        \eng{Overflow} & Hasil aritmatika terpotong; \asm{RET}
        melompat ke alamat aneh & Ketika dijalankan &
        Periksa \flag{CF}, \flag{OF}, dan pasangan \eng{dump}
        \eng{stack} \\
        \bottomrule
    \end{tabular}
\end{table}

\begin{peringatanbox}
Kesalahan sintaks dan kesalahan logika sering tertukar oleh pemula.
Perbedaannya sederhana: kesalahan sintaks membuat program
\emph{gagal dirakit}, sedangkan kesalahan logika membuat program
\emph{berhasil dirakit dan berjalan} tetapi menghasilkan nilai yang
salah. Bila assembler tidak melaporkan apa pun, jangan mencari
kesalahan pada penulisan; carilah pada nilai.
\end{peringatanbox}

Tabel~\ref{tab:taksonomi-kesalahan} menunjukkan bahwa hanya kesalahan
sintaks yang dapat dipercayakan kepada assembler. Empat jenis
kesalahan lainnya menuntut pengamatan langsung terhadap keadaan
prosesor saat program berjalan. Di sinilah \eng{debugger} menjadi
alat yang tidak tergantikan.

\subsection{Alat Debugging yang Tersedia}

Terdapat tiga alat \eng{debugging} yang dipakai pada mata kuliah ini.
Ketiganya mampu melakukan hal yang sama, yaitu menghentikan program,
memeriksa register, dan memeriksa memori; perbedaannya terletak pada
kemudahan pemakaian dan kedekatan dengan mesin sesungguhnya.

\begin{description}
    \item[MS-DEBUG di DOSBox] Program bawaan DOS yang seluruh
          interaksinya berupa baris perintah. Yang terlihat
          benar-benar isi memori, bukan tampilan yang disederhanakan
          \cite{irvine1999}.

    \item[\eng{Debugger} terpadu GUI Turbo Assembler] Satu jendela
          terpadu memuat kode sumber, register, dan memori sekaligus;
          program dijalankan baris per baris dengan menekan tombol
          \cite{guitasm}.

    \item[Jendela \eng{debug} Emu8086] Register, \eng{flag}, memori,
          dan \eng{stack} tampil sebagai panel terpisah yang selalu
          terlihat sehingga cocok untuk pemula. Karena Emu8086 adalah
          emulator, program yang dijalankan adalah hasil emulasi
          \cite{emu8086help}.
\end{description}

\begin{table}[H]
    \centering
    \caption{Perbandingan kemampuan tiga alat \eng{debugging}.}
    \label{tab:alat-debugging}
    \small
    \begin{tabular}{@{}p{2.9cm}p{3.5cm}p{3.7cm}p{3.2cm}@{}}
        \toprule
        \textbf{Kemampuan} & \textbf{MS-DEBUG (DOSBox)} &
        \textbf{GUI Turbo Assembler} & \textbf{Emu8086} \\
        \midrule
        Tampilan & Teks, baris perintah & Grafis, jendela terpadu &
        Grafis, panel terpisah \\
        \addlinespace
        Langkah per instruksi & \code{T} dan \code{P} &
        Tombol \file{Step} & Tombol \file{Single Step} \\
        \addlinespace
        \eng{Breakpoint} & \code{G} pada alamat; jumlah terbatas &
        Titik henti pada baris sumber & Titik henti pada baris sumber \\
        \addlinespace
        Memeriksa memori & \code{D} dan \code{E} &
        Jendela \eng{dump} memori & Jendela memori \\
        \addlinespace
        Memeriksa \eng{stack} & \code{D SS:SP} &
        Jendela \eng{stack} & Jendela \eng{stack} \\
        \addlinespace
        Memodifikasi kode mesin & \code{A} dan \code{E} &
        Tidak langsung & Tidak langsung \\
        \addlinespace
        Menjalankan \file{.exe} lama & Ya & Ya & Perlu emulasi \\
        \bottomrule
    \end{tabular}
\end{table}

\begin{infobox}
Pada bab ini penekanan diberikan kepada MS-DEBUG karena dua alasan.
Pertama, DEBUG memperlihatkan keadaan mesin apa adanya sehingga
mahasiswa belajar membaca alamat, opcode, dan \eng{dump} memori secara
mentah. Kedua, kemampuan membaca keluaran DEBUG berlaku umum: jendela
grafis Emu8086 maupun GUI Turbo Assembler menampilkan informasi yang
sama dengan tata letak yang lebih nyaman.
\end{infobox}
